\chapter{Green（格林/绿）函数法}

\problem{18.1}{两个相距为 $a$ 均接地的理想导体平面，两平行导体平面间放置一电荷量为 $q$ 的点电荷，电荷距离下平面为 $b,\,b<a$，求导体间的电势分布。（提示：用镜像法求解。）}
\answerarea

\problem{18.2}{一半径为 $a$ 的导体球壳接地，在球外距球心为 $b$ 处（$b>a$）放置一点电荷 $q$，试用镜像法求电场的分布，并求出导体球壳表面感应电荷的分布，进而证明此电荷分布等价于球内的镜像电荷。}
\answerarea
\clearpage

\problem{18.3}{用 Green 函数法求解第一类边值问题的球内 Laplace 方程的解。
\[
  \left\{\begin{aligned}
    &\nabla^2 u = 0,\quad r<a \\
    &u\big|_{r=a} = \cos\theta
  \end{aligned}\right.
\]}
\answerarea

\problem{18.4}{半径为 $b$ 的无穷长均匀圆柱体，初始温度为 $U_0$，初始时刻外界空间的温度为零，用 Green 函数法求解 $t$ 时刻柱外空间的温度 $u(\rho,t),\,\rho>b$。}
\answerarea
\clearpage

\problem{18.5}{侧面为绝热的无穷长细杆，初始温度为零，在 $t=0$ 时刻在 $x=x_0$ 点有一瞬时热源 $Q$，试构造该一维问题的 Green 函数。}
\answerarea
\clearpage
